\chapter{Valued networks}\label{ch:valued}

With \opt{family()} the values of the ties are modelled. The families and their parameters are given in Chapter~\ref{ch:model}. \texttt{poisson} needs non-negative integer values (counts); \texttt{binomial} needs integers from 0 to the number of trials $m_{ij}$, which \opt{trials()} gives as a positive integer (the same for all dyads), a network of the same actors or a Stata matrix; \texttt{normal} takes any real values and a dispersion $\tau^2$ with a scaled inverse chi-squared prior (\opt{dispvar()}, default 1, and \opt{dispdf()}, default 3, degrees of freedom; \latentnet{} requires the prior of the dispersion to be given, here the defaults are those). The values of the network are used as they are: a dyad without a tie has the value 0, and a missing dyad is left out of the likelihood. Sender, receiver and sociality effects, covariates, clusters, the bilinear term, chains and threads work as in the binary model, with the same priors for the coefficients, the positions and the random effects.

The start positions come from classical scaling of a dissimilarity made from the values and the coefficients from a generalised linear model of the family given those distances. The dispersion of the normal family is sampled by a Gibbs step and kept as the last row of \texttt{e(vsummary)} (and as the variable \texttt{dispersion} of \opt{savedraws()}). With a valued family \texttt{e(P)} is the matrix of the posterior mean expected values $E(y_{ij})$ (the tie probabilities for the binary model, $m_{ij}$ times the probability for the binomial family).

\section{An example: the liking of the monks}

In Sampson's data \texttt{samplike} the liking between the monks is a value from 0 to 3 (the network has 88 ties with a positive value; a dyad without a value has 0). We treat the values as the number of successes out of 3 trials in a binomial model with three clusters, and look at the fit with \texttt{estat gof}, once with the default cutoff (every positive value is a tie) and once with the cutoff 1, and at \texttt{estat ppcheck}:
\showlog{monks_valued}
\paragraph{Goodness of fit for valued networks.} \texttt{estat gof} works on the network of the dyads whose value exceeds \opt{cutoff()}, by default 0 for the Poisson and binomial families. The statistic \texttt{values} tabulates the number of dyads with each value. After \opt{family(normal)} the structure statistics are computed only if \opt{cutoff()} is given. The observed and the simulated networks are cut in the same way.

The example shows the syntax and the output, and it shows a model that does not fit. The \texttt{values} table has observed counts of 38, 20 and 30 dyads with the values 1, 2 and 3, against simulated means of 62.4, 31.9 and 13.3 (Monte Carlo $p$-values 0.00 in all three rows): the binomial model with a common number of trials spreads the positive values over the dyads too evenly and has too few dyads with the value 3. \texttt{estat ppcheck} agrees: the root mean squared error of the posterior mean expected values is 0.69 and their correlation with the observed values 0.72. With the cutoff 1 the distance statistic is reproduced (no Monte Carlo $p$-value below 0.08). A model that fits these data better would need to treat the values as ordered categories, which \nwlsm{} does not offer (Chapter~\ref{ch:limits}).
